% GENERATED FILE. Edit canonical lessons, then run npm run generate:books.
% canonical-source-hash: fnv1a64:295c2b96c86e4888
% canonical-lessons: MR-C251-prayatna-karne, MR-C251-calu-thevne, MR-C251-thambavne, MR-C251-halavne, MR-C251-garam-karne

\chapter{To try, To keep going, To bring to a stop, To move, To heat}
\label{ch:mr-to-try-to-keep-going-to-bring-to-a-stop-to-move-to-heat}
\hlchaptermodality{251}

\begin{chapteropening}
\textbf{By the end of this chapter:} \emph{I can say to try, to keep going, to bring to a stop, to move and to heat.}

Complete the last lesson of chapter 251: I can say to try, to keep going, to bring to a stop, to move and to heat.
\end{chapteropening}

\section[\texorpdfstring{\mr{प्रयत्न} \mr{करणे} (prayatna karṇe)}{प्रयत्न करणे (prayatna karṇe)}]{\mr{प्रयत्न} \mr{करणे} (prayatna karṇe) — to try}

\label{lesson:MR-C251-prayatna-karne}



\begin{quote}
Before the new one: say the Marathi for art, then the Marathi for a team.
\end{quote}

\subsection*{You'll want to know: \mr{प्रयत्न} \mr{करणे}}
\textbf{\mr{प्रयत्न} \mr{करणे}} — \emph{prayatna karṇe} — \textquotedblleft{}to try\textquotedblright{}.

\begin{grammarlens}[title={What you've built}]
One more word for this chapter.
\end{grammarlens}

\subsection*{Your turn}
\begin{itemize}
\raggedright
  \item \emph{Say it:} \emph{prayatna karṇe}
  \item \emph{Say it:} \emph{prayatna karṇe}, once more
  \item \emph{Say it:} say \emph{prayatna karṇe}
  \item \emph{Recall:} say the Marathi for art, then the Marathi for a team, then say \emph{prayatna karṇe} again
\end{itemize}

\begin{tcolorbox}[breakable,colback=teal!4,colframe=teal!35!black,title={Before you move on}]
What does \mr{प्रयत्न} \mr{करणे} mean? (\textquotedblleft{}To try\textquotedblright{}.) Say it once more.
\end{tcolorbox}
\section[\texorpdfstring{\mr{चालू} \mr{ठेवणे} (cālū ṭhevṇe)}{चालू ठेवणे (cālū ṭhevṇe)}]{\mr{चालू} \mr{ठेवणे} (cālū ṭhevṇe) — to keep going, to continue}

\label{lesson:MR-C251-calu-thevne}



\begin{quote}
Before the new one: say the Marathi for a team, then the Marathi for to try.
\end{quote}

\subsection*{You'll want to know: \mr{चालू} \mr{ठेवणे}}
\textbf{\mr{चालू} \mr{ठेवणे}} — \emph{cālū ṭhevṇe} — \textquotedblleft{}to keep going, to continue\textquotedblright{}.

\begin{grammarlens}[title={What you've built}]
One more word for this chapter.
\end{grammarlens}

\subsection*{Your turn}
\begin{itemize}
\raggedright
  \item \emph{Say it:} \emph{cālū ṭhevṇe}
  \item \emph{Say it:} \emph{cālū ṭhevṇe}, once more
  \item \emph{Say it:} say \emph{cālū ṭhevṇe}
  \item \emph{Recall:} say the Marathi for a team, then the Marathi for to try, then say \emph{cālū ṭhevṇe} again
\end{itemize}

\begin{tcolorbox}[breakable,colback=teal!4,colframe=teal!35!black,title={Before you move on}]
What does \mr{चालू} \mr{ठेवणे} mean? (\textquotedblleft{}To keep going, to continue\textquotedblright{}.) Say it once more.
\end{tcolorbox}
\section[\texorpdfstring{\mr{थांबवणे} (thāmbavṇe)}{थांबवणे (thāmbavṇe)}]{\mr{थांबवणे} (thāmbavṇe) — to bring to a stop}

\label{lesson:MR-C251-thambavne}



\begin{quote}
Before the new one: say the Marathi for to try, then the Marathi for to keep going.
\end{quote}

\subsection*{You'll want to know: \mr{थांबवणे}}
\textbf{\mr{थांबवणे}} — \emph{thāmbavṇe} — \textquotedblleft{}to bring to a stop\textquotedblright{}.

\begin{grammarlens}[title={What you've built}]
One more word for this chapter.
\end{grammarlens}

\subsection*{Your turn}
\begin{itemize}
\raggedright
  \item \emph{Say it:} \emph{thāmbavṇe}
  \item \emph{Say it:} \emph{thāmbavṇe}, once more
  \item \emph{Say it:} say \emph{thāmbavṇe}
  \item \emph{Recall:} say the Marathi for to try, then the Marathi for to keep going, then say \emph{thāmbavṇe} again
\end{itemize}

\begin{tcolorbox}[breakable,colback=teal!4,colframe=teal!35!black,title={Before you move on}]
What does \mr{थांबवणे} mean? (\textquotedblleft{}To bring to a stop\textquotedblright{}.) Say it once more.
\end{tcolorbox}
\section[\texorpdfstring{\mr{हलवणे} (halavṇe)}{हलवणे (halavṇe)}]{\mr{हलवणे} (halavṇe) — to move (something)}

\label{lesson:MR-C251-halavne}



\begin{quote}
Before the new one: say the Marathi for to keep going, then the Marathi for to bring to a stop.
\end{quote}

\subsection*{You'll want to know: \mr{हलवणे}}
\textbf{\mr{हलवणे}} — \emph{halavṇe} — \textquotedblleft{}to move (something)\textquotedblright{}.

\begin{grammarlens}[title={What you've built}]
One more word for this chapter.
\end{grammarlens}

\subsection*{Your turn}
\begin{itemize}
\raggedright
  \item \emph{Say it:} \emph{halavṇe}
  \item \emph{Say it:} \emph{halavṇe}, once more
  \item \emph{Say it:} say \emph{halavṇe}
  \item \emph{Recall:} say the Marathi for to keep going, then the Marathi for to bring to a stop, then say \emph{halavṇe} again
\end{itemize}

\begin{tcolorbox}[breakable,colback=teal!4,colframe=teal!35!black,title={Before you move on}]
What does \mr{हलवणे} mean? (\textquotedblleft{}To move (something)\textquotedblright{}.) Say it once more.
\end{tcolorbox}
\section[\texorpdfstring{\mr{गरम} \mr{करणे} (garam karṇe)}{गरम करणे (garam karṇe)}]{\mr{गरम} \mr{करणे} (garam karṇe) — to heat}

\label{lesson:MR-C251-garam-karne}



\begin{quote}
Before the new one: say the Marathi for to bring to a stop, then the Marathi for to move.
\end{quote}

\subsection*{You'll want to know: \mr{गरम} \mr{करणे}}
\textbf{\mr{गरम} \mr{करणे}} — \emph{garam karṇe} — \textquotedblleft{}to heat\textquotedblright{}.

\begin{grammarlens}[title={What you've built}]
One more word for this chapter.
\end{grammarlens}

\subsection*{Your turn}
\begin{itemize}
\raggedright
  \item \emph{Say it:} \emph{garam karṇe}
  \item \emph{Say it:} \emph{garam karṇe}, once more
  \item \emph{Say it:} say \emph{garam karṇe}
  \item \emph{Recall:} say the Marathi for to bring to a stop, then the Marathi for to move, then say \emph{garam karṇe} again
\end{itemize}

\begin{tcolorbox}[breakable,colback=teal!4,colframe=teal!35!black,title={Before you move on}]
What does \mr{गरम} \mr{करणे} mean? (\textquotedblleft{}To heat\textquotedblright{}.) Say it once more.
\end{tcolorbox}
